\chapter*{\xiaosan\heiti{摘\quad 要}}
\addcontentsline{toc}{chapter}{摘要}

数据库作为计算机系统的核心组件，其可靠性与性能始终是研究和实践的关键目标。分析型数据库因其需处理多源异构数据并支持复杂分析查询，对查询效率提出了更高要求。为突破性能瓶颈，本文聚焦于数据库查询处理流程（涵盖解析器、规划器、优化器、执行器等阶段），提出了一种基于布隆过滤器和机器学习的动态缓存加速创新方法以显著提升查询速度。本文的主要贡献如下：

第一，基于布隆过滤器的SQL/CTE动态缓存策略设计：本文深入剖析了数据库查询处理的各个阶段，识别出重复或相似查询结果复用不足的瓶颈。针对此，创新性地提出了基于SQL查询结构及公共表表达式（Common Table Expression, CTE）的动态缓存策略。该策略结合布隆过滤器技术，实现了对可复用缓存结果的快速识别与精准匹配，从而绕过部分处理流程，直接返回缓存结果，大幅降低查询延迟。针对不同查询模式，本文进行了详细的实验验证与性能分析。

第二，智能缓存替换策略：为提升缓存命中率并确保缓存数据的时效性，本文设计并实现了多策略融合的智能缓存替换系统，包括经典的LRU策略、基于时间窗口的TTL策略、混合策略以及基于机器学习的智能更新策略，可根据不同应用场景动态选择最优策略。根据配置系统通过策略协调机制实现运行时动态切换，通过自适应优化机制根据性能反馈持续优化缓存策略，通过多策略并行评估和智能决策融合实现高效而准确的淘汰决策。

第三，多策略融合的持久化存储技术：为确保缓存安全性与故障恢复能力，本文提出了四种持久化方案：内存模式、WAL模式、物化视图模式和混合模式，根据数据特征和系统状态动态选择最优持久化策略。同时，创新性地开发了基于文件系统的跨进程缓存共享机制，通过内存映射技术和一致性保证机制，实现了高效的跨进程数据共享，有效防止缓存数据丢失。

本文基于开源分析型数据库DuckDB v1.3.2进行开发，采用TPC-H、TPC-DS标准基准测试和多种自定义数据集，对缓存系统进行了系统性验证。在TPC标准基准测试中，系统实现了平均54.2\%的性能提升；在专项功能测试中获得平均68.6\%的性能提升，最高达88.2\%；布隆过滤器优化实现最高14.5\%的查询加速；创新性的文件系统跨进程缓存机制实现了最高264倍的平均加速比和99.61\%的性能提升。测试覆盖了从简单聚合到复杂CTE查询的场景，验证了缓存技术在不同查询复杂度下的稳定优化效果。

\vspace{0.5cm}
% \hspace{-1cm}
\sihao{\heiti{关键词：}}\xiaosi{数据库系统，查询缓存、CTE、缓存策略，缓存持久化}
